\section{Notation and conventions}

\begin{center}
\begin{tabular}{ll}
\toprule
Symbol & Meaning\\
\midrule
$\mathcal S$ & complete continuous physical scene\\
$\mathcal G$ & geometry and object poses\\
$\mathcal M$ & material constitutive laws\\
$P$ & number of electrical ports\\
$i,v$ & port current and voltage phasors/vectors\\
$Z$ & impedance or transfer matrix\\
$x_{\mathrm{EM}}$ & integral-equation electromagnetic coefficients\\
$\mathcal A_{\mathrm{EM}}$ & physical electromagnetic integral operator\\
$D_{\mathrm{abs}}$ & integrated absorbed-power operator\\
$H_q(\xi)$ & continuous local Hermitian loss operator\\
$L_q(\xi)$ & PSD factor such that $H_q=L_q\adj L_q$\\
$p(t)$ & low-dimensional thermal source-channel amplitudes\\
$z(t)$ & reduced thermal state\\
$C_T,G_T$ & reduced heat-capacity and dissipation matrices\\
$\Phi_T(\mathbf r)$ & continuous temperature decoder\\
$\eta_{\mathrm{EM}}$ & true normalized electromagnetic residual\\
\bottomrule
\end{tabular}
\end{center}

The document uses the engineering phasor convention consistently within each
derived formulation; an implementation must declare whether time dependence is
$e^{+j\omega t}$ or $e^{-j\omega t}$ and adjust the signs of imaginary
constitutive losses accordingly. The theory relies on physical passivity, not
on one particular sign convention.

A superscript $\adj$ denotes conjugate transpose, while $\T$ denotes ordinary
transpose. $\Herm A=(A+A\adj)/2$. Positive-semidefinite and
positive-definite matrices are denoted by $\PSD$ and $\SPD$, respectively.
